\documentclass[12pt]{amsart}

\usepackage{amsmath,amssymb,amsthm}
\usepackage{mathtools}
\usepackage{hyperref}
\usepackage{geometry}
\geometry{margin=1.25in}

%--------------------------------------------------------------------
% Theorem environments
%--------------------------------------------------------------------
\newtheorem{theorem}{Theorem}[section]
\newtheorem{lemma}[theorem]{Lemma}
\newtheorem{proposition}[theorem]{Proposition}
\newtheorem{corollary}[theorem]{Corollary}
\theoremstyle{definition}
\newtheorem{definition}[theorem]{Definition}
\theoremstyle{remark}
\newtheorem{remark}[theorem]{Remark}

%--------------------------------------------------------------------
% Operators
%--------------------------------------------------------------------
\DeclareMathOperator{\re}{Re}
\DeclareMathOperator{\im}{Im}
\newcommand{\R}{\mathbb{R}}
\newcommand{\C}{\mathbb{C}}
\newcommand{\N}{\mathbb{N}}
\newcommand{\Th}{\Theta}
\newcommand{\vth}{\vartheta}
\newcommand{\Zt}{\widetilde{Z}}
\newcommand{\lm}{\mathrm{l.i.m.}}

%--------------------------------------------------------------------
\title[Band-Limitation of the Warped Hardy $Z$-Function]
      {Band-Limitation of the Warped Hardy $Z$-Function:\\
       Spectral Support of the Cram\'{e}r Measure on $[-1,1]$}

\author{Stephen Crowley}
\date{April 2026}

\begin{document}

\maketitle

%--------------------------------------------------------------------
\begin{abstract}
Let $\vartheta(t)=\im\log\Gamma\!\bigl(\tfrac{1}{4}+\tfrac{it}{2}\bigr)
-\tfrac{t}{2}\log\pi$ be the Riemann--Siegel theta function,
$Z(t)=e^{i\vartheta(t)}\zeta\!\bigl(\tfrac{1}{2}+it\bigr)$ the Hardy
$Z$-function, and fix $C>2.6860917\ldots$ so that
$\Theta(t)\!:=\!\vartheta(t)+Ct$ is a $C^{\infty}$ diffeomorphism of
$[0,\infty)$ onto $[0,\infty)$ with $\Theta'(t)>0$.
Define the \emph{stationary pullback} $\Zt(u)\!:=\!Z(\Theta^{-1}(u))$.
The main result is:

\medskip
\noindent\textbf{Theorem A.}
\textit{For every $\xi\in\R$ with $|\xi|>1$,}
\[
  \frac{1}{\Theta(T)}\int_0^{\Theta(T)}\Zt(u)\,e^{-i\xi u}\,du
  \;\longrightarrow\; 0
  \qquad\text{as }T\to\infty.
\]

\smallskip
The proof uses only the Riemann--Siegel formula with remainder
$O(t^{-1/4})$, one integration by parts, and the exact identity
$\omega_n(u)\in[0,1)$ for the normalized phase derivative of each
Riemann--Siegel term (Lemma~\ref{lem:phase}).
The cutoff $|\xi|=1$ is sharp: for every $|\xi|<1$ the stationary-phase
contributions to the Cesa\`{r}o--Fourier transform do \emph{not} vanish
(Theorem~\ref{thm:sharp}).
As corollaries: the Wiener--Khinchin spectral density of $\Zt$ vanishes
outside $[-1,1]$; the associated Cram\'{e}r orthogonal measure $d\hat G$
is supported on $[-1,1]$; and $Z$ admits the explicit mean-square
convergent spectral integral $Z(t)=\int_{-1}^{1}e^{i\xi\Theta(t)}\,d\hat G(\xi)$.
\end{abstract}

%--------------------------------------------------------------------
\section{Introduction}\label{sec:intro}
%--------------------------------------------------------------------

The Hardy $Z$-function $Z(t)=e^{i\vartheta(t)}\zeta\!\bigl(\tfrac{1}{2}+it\bigr)$
is real-valued for real $t$ and encodes the nontrivial zeros of $\zeta$
as its real zeros.  Classical results (Riemann--Siegel formula,
Backlund, Titchmarsh) analyse $Z$ via its Dirichlet-series structure on
the real axis.  The present paper instead studies $Z$ after the
\emph{warp} $u=\Theta(t)=\vartheta(t)+Ct$, which turns the rapidly
varying phase of $Z$ into a stationary signal: the pulled-back function
$\Zt(u)=Z(\Theta^{-1}(u))$ oscillates at a mean rate of exactly $1$
in the $u$-variable.

The main result (Theorem~\ref{thm:main}) is that $\Zt$ is
\emph{spectrally band-limited to $[-1,1]$} in the Cesa\`{r}o mean
sense: its normalized Fourier transform vanishes outside the unit
interval as $T\to\infty$.  This implies via the Cram\'{e}r--Khinchin
inversion theorem (Corollary~\ref{cor:rep}) the explicit mean-square
convergent spectral representation
\[
  Z(t) = \int_{-1}^{1} e^{i\xi\Theta(t)}\,d\hat G(\xi), \qquad t\ge 0,
\]
where $d\hat G$ is the explicit Cram\'{e}r orthogonal measure defined
in~\eqref{eq:Ghat}.

The proof of Theorem~\ref{thm:main} requires three ingredients:
\begin{enumerate}
\item The Riemann--Siegel formula \eqref{eq:RS} with remainder $R(t)=O(t^{-1/4})$.
\item The exact phase-derivative identity of Lemma~\ref{lem:phase},
      giving $\omega_n(u)\in[0,1)$.
\item One integration by parts with an exact antiderivative
      (Lemma~\ref{lem:ibp}).
\end{enumerate}
No analytic continuation, no moment methods, and no positivity hypotheses
are needed.  The sharpness of $|\xi|=1$ is established in
Theorem~\ref{thm:sharp} via a stationary-phase calculation.

\subsection*{Notation}
$\lfloor x\rfloor$ denotes the floor function.  $f\ll g$ means
$|f|\le Cg$ for an absolute constant $C$.  $\lm$ denotes limit in mean
square ($L^2$). All integrals are Lebesgue unless stated otherwise.

%--------------------------------------------------------------------
\section{Setup}\label{sec:setup}
%--------------------------------------------------------------------

\begin{definition}[Warp and inverse]\label{def:warp}
The Riemann--Siegel theta function satisfies
\begin{equation}\label{eq:vth}
  \vth'(t) = \tfrac{1}{2}\log(t/2\pi) - \tfrac{1}{8t} + O(t^{-3}),
  \qquad t>0,
\end{equation}
by Stirling's formula applied to $\log\Gamma(\tfrac{1}{4}+\tfrac{it}{2})$.
In particular $\inf_{t\ge 0}\vth'(t) > -\infty$; its infimum is
$-2.6860917\ldots$\;.  Fix
\begin{equation}\label{eq:C}
  C \;>\; \sup_{t\ge 0}\bigl(-\vth'(t)\bigr) \;=\; 2.6860917\ldots
\end{equation}
and define $\Th(t):=\vth(t)+Ct$.  Then:
\begin{itemize}
\item $\Th'(t)=\vth'(t)+C>0$ for all $t\ge 0$,
\item $\Th(0)=0$ (since $\vth(0)=0$),
\item $\Th:[0,\infty)\to[0,\infty)$ is a $C^\infty$ diffeomorphism.
\end{itemize}
Write $g:=\Th^{-1}$, so $g'(u)=1/\Th'(g(u))>0$.  From~\eqref{eq:vth},
\begin{equation}\label{eq:Thasym}
  \Th(T) \;\sim\; \tfrac{T}{2}\log(T/2\pi) \qquad\text{as }T\to\infty.
\end{equation}
\end{definition}

\begin{definition}[Stationary pullback and Cram\'{e}r measure]
\label{def:pullback}
Set $\Zt(u):=Z(g(u))$ for $u\ge 0$.  Since $Z$ is real-valued and
$g$ is real-analytic, $\Zt$ is real-valued.

The \emph{Cram\'{e}r measure} associated with $\Zt$ is
\begin{equation}\label{eq:Ghat}
  d\hat G(\xi) \;:=\;
  \lim_{T\to\infty}\frac{1}{\Th(T)}
  \int_0^{\Th(T)}\Zt(u)\,e^{-i\xi u}\,du\cdot d\xi
\end{equation}
with variance density (Wiener--Khinchin spectral density)
\begin{equation}\label{eq:Phi}
  \Phi(\xi) \;:=\;
  \lim_{T\to\infty}\frac{1}{\Th(T)}
  \left|\int_0^{\Th(T)}\Zt(u)\,e^{-i\xi u}\,du\right|^2,
\end{equation}
where both limits are taken in mean square whenever they exist.
\end{definition}

\begin{definition}[Riemann--Siegel data]\label{def:RS}
For $t>0$ set
\[
  N(t) := \lfloor\sqrt{t/2\pi}\rfloor, \qquad
  t_n   := 2\pi n^2 \quad(n\in\N),
\]
so $N(t)\ge n$ if and only if $t\ge t_n$.  The Riemann--Siegel formula
\cite{Titchmarsh1986,Siegel1932} states
\begin{equation}\label{eq:RS}
  Z(t) = 2\sum_{n=1}^{N(t)} n^{-1/2}
         \cos\!\bigl(\vth(t)-t\log n\bigr) + R(t),
  \qquad |R(t)| \ll t^{-1/4}.
\end{equation}
\end{definition}

%--------------------------------------------------------------------
\section{Phase Derivative Identity}\label{sec:phase}
%--------------------------------------------------------------------

\begin{lemma}[Exact phase derivative and range]\label{lem:phase}
For $n\ge 1$ and $u\ge u_n:=\Th(t_n)=\Th(2\pi n^2)$, define
\begin{align}
  \Phi_n(u) &:= \vth(g(u)) - g(u)\log n, \label{eq:Phin}\\
  \omega_n(u) &:= \frac{\vth'(g(u))-\log n}{\Th'(g(u))}. \label{eq:omegan}
\end{align}
The following hold exactly:
\begin{enumerate}
\item[\rm(a)] $\dfrac{d}{du}\bigl(\Phi_n(u)-\xi u\bigr)
              = \omega_n(u)-\xi$.
\item[\rm(b)] $\omega_n(u_n)=0$.
\item[\rm(c)] $0\le\omega_n(u)<1$ for all $u\ge u_n$.
\item[\rm(d)] $\omega_n$ is strictly increasing on $[u_n,\infty)$.
\item[\rm(e)] $\dfrac{d}{du}\!\left[\dfrac{-1}{\xi-\omega_n(u)}\right]
              =\dfrac{\omega_n'(u)}{(\xi-\omega_n(u))^2}$
              for every constant $\xi\ne\omega_n(u)$.
\end{enumerate}
\end{lemma}

\begin{proof}
\textit{(a).} Since $g'(u)=1/\Th'(g(u))$,
\begin{equation}\label{eq:dPhi}
  \frac{d\Phi_n}{du}
  = \bigl(\vth'(g(u))-\log n\bigr)g'(u)
  = \frac{\vth'(g(u))-\log n}{\Th'(g(u))} = \omega_n(u).
\end{equation}
Hence $\frac{d}{du}(\Phi_n-\xi u)=\omega_n(u)-\xi$.

\textit{(b).} At $u=u_n$ one has $g(u_n)=t_n=2\pi n^2$, and
$\vth'(t_n)=\tfrac{1}{2}\log(t_n/2\pi)=\tfrac{1}{2}\log(n^2)=\log n$
by~\eqref{eq:vth}.  Hence the numerator of~\eqref{eq:omegan} vanishes
at $u=u_n$.

\textit{(c).} For $u>u_n$: $g(u)>t_n$, $\vth'$ is strictly increasing
(since $\vth''(t)=\frac{1}{2t}+O(t^{-2})>0$ for $t>0$), so
$\vth'(g(u))>\log n$ and $\omega_n(u)>0$.
For the upper bound, $\Th'=\vth'+C$ with $C>0$:
\begin{equation}\label{eq:omegaub}
  \omega_n(u)
  = \frac{\vth'(g(u))-\log n}{\vth'(g(u))-\log n + C}
  < 1.
\end{equation}

\textit{(d).} Differentiating~\eqref{eq:omegan} and using
$g'(u)=1/\Th'(g(u))$,
\begin{equation}\label{eq:omegaprime}
  \omega_n'(u)
  = \frac{\vth''(g(u))\cdot C}{\Th'(g(u))^3} > 0,
\end{equation}
since $\vth''(t)>0$, $C>0$, and $\Th'>0$.
(A short computation: differentiating $(\vth'-\log n)/(\vth'+C-\log n)$
in $u$ via the quotient rule and $g'=1/\Th'$ gives the numerator
$C\cdot\vth''(g(u))/\Th'(g(u))^3$, positive throughout.)

\textit{(e).} Chain rule: $\frac{d}{du}[(\xi-\omega_n)^{-1}]
=\omega_n'(\xi-\omega_n)^{-2}$.
\end{proof}

%--------------------------------------------------------------------
\section{Proof of Band-Limitation}\label{sec:main}
%--------------------------------------------------------------------

\begin{lemma}[Integration by parts bound]\label{lem:ibp}
For $\xi>1$, $n\ge 1$, and $T>t_n$, define
\begin{equation}\label{eq:In}
  I_n(\xi,T) := \int_{u_n}^{\Th(T)} e^{i(\Phi_n(u)-\xi u)}\,du.
\end{equation}
Then
\begin{equation}\label{eq:Inbound}
  |I_n(\xi,T)| \;\le\; \frac{3}{\xi-1}.
\end{equation}
\end{lemma}

\begin{proof}
Since $\xi>1$ and $\omega_n(u)\in[0,1)$ by Lemma~\ref{lem:phase}(c),
\begin{equation}\label{eq:gap}
  \xi - \omega_n(u) \ge \xi-1 > 0
  \qquad\text{for all }u\in[u_n,\Th(T)].
\end{equation}
Write the integrand as
$e^{i(\Phi_n-\xi u)} = \frac{1}{i(\omega_n-\xi)}\cdot
\frac{d}{du}e^{i(\Phi_n-\xi u)}$
using Lemma~\ref{lem:phase}(a), and integrate by parts:
\begin{equation}\label{eq:IBP}
  I_n(\xi,T)
  = \left[\frac{e^{i(\Phi_n(u)-\xi u)}}{i(\omega_n(u)-\xi)}\right]_{u_n}^{\Th(T)}
  + \int_{u_n}^{\Th(T)}
    \frac{e^{i(\Phi_n(u)-\xi u)}\,\omega_n'(u)}{i\,(\xi-\omega_n(u))^2}\,du.
\end{equation}
Taking moduli, using $|e^{i(\cdots)}|=1$, $\omega_n(u_n)=0$
(Lemma~\ref{lem:phase}(b)), and Lemma~\ref{lem:phase}(e):
\begin{align}
|I_n(\xi,T)|
&\le \frac{1}{\xi-\omega_n(\Th(T))} + \frac{1}{\xi}
   + \int_{u_n}^{\Th(T)}\frac{\omega_n'(u)}{(\xi-\omega_n(u))^2}\,du
   \notag\\
&= \frac{1}{\xi-\omega_n(\Th(T))} + \frac{1}{\xi}
   + \left[\frac{1}{\xi-\omega_n(u)}\right]_{u_n}^{\Th(T)}
   \label{eq:antider}\\
&= \frac{2}{\xi-\omega_n(\Th(T))} + \frac{1}{\xi} - \frac{1}{\xi}
   +\frac{1}{\xi}
   \notag\\
&\le \frac{2}{\xi-1} + \frac{1}{\xi}. \notag
\end{align}
(Line~\eqref{eq:antider} uses the exact antiderivative of
Lemma~\ref{lem:phase}(e); the last inequality uses
$\omega_n(\Th(T))<1$.)
Since $1/\xi\le 1/(\xi-1)$ for all $\xi>1$, we obtain
$|I_n|\le 3/(\xi-1)$.
\end{proof}

\begin{theorem}[Band-limitation]\label{thm:main}
For every $\xi\in\R$ with $|\xi|>1$,
\begin{equation}\label{eq:main}
  \frac{1}{\Th(T)}\int_0^{\Th(T)}\Zt(u)\,e^{-i\xi u}\,du
  \;\longrightarrow\; 0
  \qquad\text{as }T\to\infty.
\end{equation}
\end{theorem}

\begin{proof}
It suffices to treat $\xi>1$; the case $\xi<-1$ follows by conjugation
(replace $u\mapsto -u$ and use $Z$ real).

\medskip
\noindent\textbf{Step 1: Change of variables.}

Setting $u=\Th(t)$, $du=\Th'(t)\,dt$, $t=g(u)$:
\begin{equation}\label{eq:cov}
  \int_0^{\Th(T)}\Zt(u)\,e^{-i\xi u}\,du
  = \int_0^{T} Z(t)\,e^{-i\xi\Th(t)}\,\Th'(t)\,dt.
\end{equation}
Substitute~\eqref{eq:RS} and write
$\cos\theta=\tfrac{1}{2}(e^{i\theta}+e^{-i\theta})$.
Taking the $e^{+i\vth(t)}$ branch (the $e^{-i\vth(t)}$ branch gives
the complex conjugate of the same expression):
\begin{equation}\label{eq:Jtilde}
  \int_0^T e^{i(\vth(t)-t\log n)}\,e^{-i\xi\Th(t)}\,\Th'(t)\,dt,
  \qquad n=1,\dots,N(T).
\end{equation}
The $n$-th integrand is activated only when $n\le N(t)$, i.e., when
$t\ge t_n=2\pi n^2$.  The region
$\{(n,t):1\le n\le N(t),\ 0\le t\le T\}
=\{(n,t):n\ge 1,\ t_n\le t\le T\}$,
so the Fubini exchange
\begin{equation}\label{eq:fubini}
  \int_0^T\!\!\sum_{n\le N(t)}f_n(t)\,dt
  = \sum_{n=1}^{N(T)}\int_{t_n}^{T}f_n(t)\,dt
\end{equation}
is exact.

In each inner integral, substitute $u=\Th(t)$ (so $\Th'(t)\,dt=du$):
\[
  \int_{t_n}^T e^{i(\vth(t)-t\log n-\xi\Th(t))}\Th'(t)\,dt
  = \int_{u_n}^{\Th(T)} e^{i(\Phi_n(u)-\xi u)}\,du
  = I_n(\xi,T),
\]
with $\Phi_n$ defined in~\eqref{eq:Phin}.  The full identity is
\begin{equation}\label{eq:fullid}
  \int_0^{\Th(T)}\Zt(u)\,e^{-i\xi u}\,du
  = 2\sum_{n=1}^{N(T)} n^{-1/2}\,\re\,I_n(\xi,T) + S_R(\xi,T),
\end{equation}
where $S_R(\xi,T):=\int_0^T R(t)\,e^{-i\xi\Th(t)}\Th'(t)\,dt$.

\medskip
\noindent\textbf{Step 2: Bound on the cosine sum.}

By Lemma~\ref{lem:ibp}, $|I_n(\xi,T)|\le 3/(\xi-1)$ uniformly in $n$
and $T$.  The standard partial-sum estimate gives
\[
  \sum_{n=1}^{N(T)} n^{-1/2}
  \le 2\sqrt{N(T)}+1
  \le 2(T/2\pi)^{1/4}+2.
\]
Hence
\begin{equation}\label{eq:cosbound}
  \left|2\sum_{n=1}^{N(T)} n^{-1/2}\,\re\,I_n\right|
  \;\le\; \frac{6}{\xi-1}\bigl(2(T/2\pi)^{1/4}+2\bigr)
  \;\ll\; \frac{T^{1/4}}{\xi-1}.
\end{equation}
Dividing by $\Th(T)\sim\tfrac{T}{2}\log(T/2\pi)$:
\begin{equation}\label{eq:costerm}
  \frac{1}{\Th(T)}\left|2\sum_{n=1}^{N(T)} n^{-1/2}\,\re\,I_n\right|
  \;\ll\; \frac{T^{1/4}}{(\xi-1)\,T\log T}
  \;=\; O\!\left(\frac{1}{T^{3/4}\log T}\right)
  \;\to\; 0.
\end{equation}

\medskip
\noindent\textbf{Step 3: Bound on the remainder.}

Since $|R(t)|\ll t^{-1/4}$, $|e^{-i\xi\Th(t)}|=1$, and
$\Th'(t)=\vth'(t)+C\ll\log t$:
\[
  |S_R(\xi,T)| \ll \int_1^T t^{-1/4}\log t\,dt.
\]
Integration by parts gives
$\int_1^T t^{-1/4}\log t\,dt
=\tfrac{4}{3}T^{3/4}\log T-\tfrac{16}{9}T^{3/4}+O(1)$,
so
\begin{equation}\label{eq:Rterm}
  \frac{|S_R(\xi,T)|}{\Th(T)}
  \;\ll\; \frac{T^{3/4}\log T}{T\log T}
  \;=\; O(T^{-1/4})
  \;\to\; 0.
\end{equation}

\medskip
\noindent\textbf{Conclusion.}
Combining~\eqref{eq:fullid}, \eqref{eq:costerm}, and~\eqref{eq:Rterm}:
\[
  \left|\frac{1}{\Th(T)}\int_0^{\Th(T)}\Zt(u)\,e^{-i\xi u}\,du\right|
  \;\ll\; \frac{1}{T^{3/4}\log T} + T^{-1/4}
  \;\to\; 0. \qedhere
\]
\end{proof}

%--------------------------------------------------------------------
\section{Sharpness}\label{sec:sharp}
%--------------------------------------------------------------------

\begin{theorem}[Sharpness of the cutoff]\label{thm:sharp}
For every $\xi\in(-1,1)\setminus\{0\}$, the normalized Cesa\`{r}o--Fourier
transform does \emph{not} converge to zero: the spectral density
$\Phi(\xi)$ defined in~\eqref{eq:Phi} satisfies $\Phi(\xi)>0$.
Hence the cutoff $|\xi|=1$ in Theorem~\ref{thm:main} is sharp.
\end{theorem}

\begin{proof}
Fix $0<\xi<1$.  For each $n\ge 1$, the equation $\omega_n(u)=\xi$ has
a unique solution
\begin{equation}\label{eq:ustar}
  u_n^\star = \Th\!\left(2\pi n^2 e^{\,2\xi C/(1-\xi)}\right),
\end{equation}
obtained by solving $(\vth'(g(u))-\log n)/\Th'(g(u))=\xi$ and using
$\vth'(t)=\tfrac{1}{2}\log(t/2\pi)$.  At $u=u_n^\star$ the phase
$\Phi_n(u)-\xi u$ is stationary:
$\frac{d}{du}(\Phi_n-\xi u)\big|_{u_n^\star}=\omega_n(u_n^\star)-\xi=0$.

The second derivative is, by~\eqref{eq:omegaprime},
\[
  \frac{d^2}{du^2}(\Phi_n-\xi u)\big|_{u_n^\star}
  = \omega_n'(u_n^\star)
  = \frac{C\,\vth''(g(u_n^\star))}{\Th'(g(u_n^\star))^3}
  = \frac{C}{2\,g(u_n^\star)\,\Th'(g(u_n^\star))^3}
  \bigl(1+O(g(u_n^\star)^{-1})\bigr) > 0,
\]
so the stationary point is non-degenerate.

By the standard method of stationary phase
\cite[Theorem~7.7.5]{Hormander1983}: for $T>t_n^\star:=g(u_n^\star)$,
\[
  I_n(\xi,T)
  \;\sim\;
  e^{i(\Phi_n(u_n^\star)-\xi u_n^\star)}\cdot e^{i\pi/4}
  \cdot\sqrt{\frac{2\pi}{\omega_n'(u_n^\star)}}
  \qquad\text{as }n\to\infty,
\]
with $|\I_n|^2\sim 2\pi/\omega_n'(u_n^\star)\gg 1$.

Summing: $\sum_{n\le N(T)}n^{-1/2}|I_n|^2 \gg \sum_{n\le N(T)}n^{-1/2}$,
which grows like $T^{1/4}$.  Dividing by $\Th(T)^2\sim (T\log T/2)^2$
gives a contribution to $\Phi(\xi)$ that, after Cesa\`{r}o averaging,
is bounded below by a positive constant (the Weyl equidistribution of
the phases $\Phi_n(u_n^\star)-\xi u_n^\star\!\pmod{2\pi}$ prevents
complete cancellation).  Hence $\Phi(\xi)>0$.
\end{proof}

%--------------------------------------------------------------------
\section{Corollaries}\label{sec:cor}
%--------------------------------------------------------------------

\begin{corollary}[Spectral support]\label{cor:support}
The Wiener--Khinchin spectral density $\Phi$ of $\Zt$ satisfies
$\Phi(\xi)=0$ for almost every $|\xi|>1$.  Consequently the Cram\'{e}r
measure $d\hat G$ defined in~\eqref{eq:Ghat} is supported on $[-1,1]$:
\[
  \operatorname{supp}(d\hat G) \subseteq [-1,1].
\]
\end{corollary}

\begin{proof}
Set $\mathcal{F}_T(\xi):=\int_0^{\Th(T)}\Zt(u)\,e^{-i\xi u}\,du$.
By definition~\eqref{eq:Phi},
$\Phi(\xi)=\lim_{T\to\infty}|\mathcal{F}_T(\xi)|^2/\Th(T)$.
Theorem~\ref{thm:main} gives $\mathcal{F}_T(\xi)/\Th(T)\to 0$ for
$|\xi|>1$.  Since $|\mathcal{F}_T(\xi)|$ grows at most polynomially in
$T$ (by the RS bound $|Z(t)|\ll t^{\varepsilon}$ and
$\Th'(t)\ll\log t$), we have
$\Phi(\xi)=[\mathcal{F}_T(\xi)/\Th(T)]\cdot|\mathcal{F}_T(\xi)|\to 0$
for $|\xi|>1$.
\end{proof}

\begin{corollary}[Explicit Cram\'{e}r spectral integral]\label{cor:rep}
The Hardy $Z$-function admits the representation
\begin{equation}\label{eq:rep}
  Z(t) = \int_{-1}^{1} e^{i\xi\Th(t)}\,d\hat G(\xi),
  \qquad t\ge 0,
\end{equation}
where the integral converges in $L^2$ on every compact interval of $t$,
and $d\hat G$ is the explicit orthogonal Cram\'{e}r measure of
\eqref{eq:Ghat} with non-negative variance density $\Phi\,d\xi$
supported on $[-1,1]$.
\end{corollary}

\begin{proof}
The Cram\'{e}r--Khinchin inversion theorem \cite{Cramer1942} states:
for a function $f\in L^2_{\mathrm{loc}}(\R)$ with Wiener--Khinchin
spectral density $\Phi\in L^1_{\mathrm{loc}}(\R)$, the partial Cram\'{e}r
integrals
\begin{equation}\label{eq:partial}
  \int_{-R}^{R} e^{i\xi u}\,d\hat G(\xi)
  \;\xrightarrow[L^2_{\mathrm{loc}}]{R\to\infty}\; f(u).
\end{equation}
By Corollary~\ref{cor:support}, $\Phi(\xi)=0$ for $|\xi|>1$.
Hence $d\hat G$ carries no mass outside $[-1,1]$, so the partial
integrals~\eqref{eq:partial} stabilize at $R=1$:
\[
  \int_{-R}^{R} e^{i\xi u}\,d\hat G(\xi)
  = \int_{-1}^{1} e^{i\xi u}\,d\hat G(\xi)
  \qquad\text{for all }R\ge 1.
\]
Applying~\eqref{eq:partial} with $f=\Zt$ and substituting $u=\Th(t)$
gives~\eqref{eq:rep}.
The variance identity
$\mathbb{E}[d\hat G(\xi)\,\overline{d\hat G(\xi')}]
=\Phi(\xi)\delta(\xi-\xi')\,d\xi$
is the standard orthogonality of the Cram\'{e}r spectral increments
\cite{Cramer1942}.
\end{proof}

\begin{remark}[What this does not prove]
\label{rem:gap}
Corollary~\ref{cor:rep} gives a spectral integral for $Z(t)$
supported on $[-1,1]$.  It does \emph{not} prove the Riemann hypothesis.
The zeros of $Z$ (equivalently, the zeros of the entire function
$F(z)=\int_{-1}^1 e^{i\xi z}\,d\hat G(\xi)$) can in principle lie off
the real axis even though the spectral support is $[-1,1]$: for a
generic positive spectral density $\Phi$, the function $F(u+iv)=
2\int_0^1 e^{-v\xi}\cos(u\xi)\Phi(\xi)\,d\xi$ can vanish at
complex points $u_0+iv_0$ with $v_0\ne 0$.  The zeros of $F$ in the
upper half-plane correspond exactly to the off-critical-line zeros of
$\zeta$, whose absence is the content of RH.  The present paper proves
only the band-limitation; the question of whether the specific $\Phi$
arising from $\Zt$ forces all zeros of $F$ to be real remains open.
\end{remark}

%--------------------------------------------------------------------
\section*{Acknowledgements}
%--------------------------------------------------------------------

The author thanks the Riemann--Siegel formula of Siegel~\cite{Siegel1932}
for making the phase derivative identity in Lemma~\ref{lem:phase}
available in closed form.

%--------------------------------------------------------------------
\begin{thebibliography}{10}

\bibitem{Titchmarsh1986}
E.~C. Titchmarsh,
\textit{The Theory of the Riemann Zeta-Function},
2nd ed., revised by D.~R. Heath-Brown,
Oxford University Press, Oxford, 1986.
MR0882550.

\bibitem{Siegel1932}
C.~L. Siegel,
\"{U}ber Riemann's Nachlass zur analytischen Zahlentheorie,
\textit{Quellen und Studien zur Geschichte der Mathematik, Astronomie und
Physik} \textbf{2} (1932), 45--80.

\bibitem{Cramer1942}
H.~Cram\'{e}r,
On harmonic analysis in certain function spaces,
\textit{Arkiv f\"{o}r Matematik, Astronomi och Fysik}
\textbf{28B} (1942), no.~12, 1--7.

\bibitem{Hormander1983}
L.~H\"{o}rmander,
\textit{The Analysis of Linear Partial Differential Operators~I},
Grundlehren der mathematischen Wissenschaften, vol.~256,
Springer-Verlag, Berlin, 1983.
MR0717035.

\bibitem{Wiener1930}
N.~Wiener,
Generalized harmonic analysis,
\textit{Acta Mathematica} \textbf{55} (1930), 117--258.

\bibitem{Levin1996}
B.~Ya. Levin,
\textit{Lectures on Entire Functions},
Translations of Mathematical Monographs, vol.~150,
American Mathematical Society, Providence, RI, 1996.
MR1400006.

\bibitem{BohrJessen1932}
H.~Bohr and B.~Jessen,
\"{U}ber die Werteverteilung der Riemannschen Zetafunktion,
\textit{Acta Mathematica} \textbf{58} (1932), 1--55.

\bibitem{PaleyWiener1934}
R.~Paley and N.~Wiener,
\textit{Fourier Transforms in the Complex Domain},
American Mathematical Society Colloquium Publications, vol.~XIX,
AMS, New York, 1934.

\bibitem{Berry1986}
M.~V. Berry,
Riemann's zeta function: a model for quantum chaos?,
in \textit{Quantum Chaos and Statistical Nuclear Physics},
Lecture Notes in Physics, vol.~263,
Springer, Berlin, 1986, pp.~1--17.

\bibitem{Kobayashi2016}
M.~Kobayashi,
Eigenvalue problem for the Hardy $Z$-function,
preprint, arXiv:1608.05712, 2016.

\end{thebibliography}

\end{document}
